\section{Four Coke bottles on a table (Geometry)}
Arrange four Coke bottles on a table so that the distances between each pair of caps will be the same. Assume that
you can make a Coke bottle stand on it's head (this is a pretty strong hint!, maybe I should remove it?!?). The length
or height of each bottle is ${h}$.

\subsection{Solution}
There are two solutions:

The first solution is to put the bottles in a square where the side of the square is the length of the bottle. Bottles
next to each other will stand \emph{in reverse} (if the first is on it's head the the second will stand straight).

\begin{center}
\input{out/src/sk/bottles.tex}
\end{center}

The shape created by the bottle caps is, of course, a perfect triangular pyramid with a side of length $d=\sqrt{2h^2}$.

The second solution is to put three bottles on the vertices of an equilateral triangle standing one way. The fourth will be placed
in the middle of the triangle and will be standing reversed to the others. The problem here is to calculate the side of the triangle.
Let's use the following diagram:

\begin{center}
\input{out/src/sk/perfect_triangular_pyramid.tex}
\end{center}

What we are trying to find is $d$ as a function of $h$. Since $H$, $d/2$ and $d$ create a right angled triangle, it follows
that $H^2+(d/2)^2=d^2$ and so $H=\sqrt{3}/2d$. Since the size of the inner triangle could be computed two ways then $hH=dH'$ and
since $H'=\sqrt{H^2-(d/2)^2}$ it follows after some math that $d=\sqrt{3/2}h$ which makes $d$ about $1.224h$.

The solution built on this analysis will look like this:

\begin{center}
\input{out/src/sk/bottles2.tex}
\end{center}

As is evident the key to solving this question is the perfect triangular pyramid as both solutions are based on it and it is the only
way to put 4 elements in space being spaced equally apart. In theory you could place the bottles at any height and so
infinite solutions to this riddle could be produced by making the side of the perfect pyramid as small or as large as you want. But we are
constrained by the table and so the heights at which the caps are will be either $0$ or $h$. The two solutions above seem to be primary
examples under this constraint. Minor variations on the above two solutions could be produced by making the bottles whose caps are at table
level lay instead of stand on their heads. The first solution is more attractive than the second because it could be performed in practice
without relying on measurement tools in order to place the bottles at the calculated distance from each other as per the second solution.
